\section{Background}
\label{sec:background}

\subsection{fNIRS decoding and its constraints}
fNIRS records changes in oxygenated (HbO) and deoxygenated (HbR) haemoglobin concentration from
near-infrared light passing between sources and detectors on the scalp. The signal is
haemodynamic, slow (peak responses 5--7 s after onset), and contaminated by superficial
physiology---cardiac pulsation, respiration, and Mayer waves---which is measured directly by
\emph{short-separation channels}, i.e.\ source--detector pairs with separations of roughly
8--15 mm that sample scalp and skull rather than cortex.

Two constraints dominate decoding. First, spatial resolution: a recent analysis of the resolution
required for image reconstruction reports that at 2 cm---the continuous-wave fNIRS regime---
300-way retrieval falls to about 20\,\%, against 93\,\% for full-resolution fMRI
\citep{adamic2024resolution}. Second, probe density: single-trial decoding accuracy improves
systematically with optode density \citep{fischer2026density}. Sparse montages such as the one we
audit here (33 source--detector pairs) sit at the low end of that curve.

\subsection{Foundation models across neuroimaging}
Pretrained representation models now exist for several modalities. For EEG: LaBraM
\citep{bi2024labram}, CBraMod \citep{wang2025cbramod}, EEGPT \citep{wang2024eegpt}, and unified
EEG/MEG models such as BrainOmni \citep{brainomni2025}. For fNIRS, MtReNIRS
\citep{chae2026mtrenirs} is a 1.23M-parameter spatial--temporal transformer pretrained on twelve
public datasets (1{,}387 runs, 203{,}420 ten-second windows, exceeding 567 hours) using one
32\,GB V100, producing a 128-dimensional embedding per channel and 10-second window and accepting
arbitrary montages through standard 10-05 channel coordinates. GeoNIRS \citep{geonirs2026} takes
a geometry-aware approach, and a curated 67-dataset corpus \citep{sposso2026corpus} now supports
further pretraining.

A parallel line of work has developed architectures tailored to fNIRS physiology and montage
geometry rather than adapting generic backbones: transformer and general vision frameworks for
fNIRS classification \citep{wang2022transformer,wang2022vision}, representations that model the
delayed haemodynamic response explicitly \citep{wang2023rethinking}, a unified framework informed
by local brain activation patterns \citep{wang2026unified}, and a domain-knowledge-fused
state-space model \citep{wang2026dkssm}. That work establishes that architecture and
physiological inductive bias matter for fNIRS decoding. It is also the premise our random-encoder
control tests: if architectural inductive bias is what carries the performance, a randomly
initialised encoder of the same architecture should recover much of it---and, in our results,
it does.

\subsection{The evaluation gap}
Evaluation practice has not kept pace. For EEG, independent benchmarks now evaluate dozens of
models across many tasks and datasets. For fNIRS we could find no independent benchmark: the
published evaluations of both released models are by their own authors, and much of the
evaluation data overlaps the pretraining corpora, which makes apparent transfer gains difficult
to attribute. In the EEG literature the analogous concern is already empirically grounded:
large-scale pretraining has been reported to yield no consistent advantage over compact
convolutional baselines on some downstream tasks, and the benefit of pretraining has been found
to depend on corpus scale and on the specific benchmark.

The control we adopt here is standard elsewhere but absent from fNIRS foundation-model papers.
In vision and language, a randomly initialised model of the same architecture is the accepted
baseline that separates \emph{architecture} from \emph{pretraining}; the idea has an established
lineage in the random-features literature \citep{rahimi2007random}, where fixed random
projections are known to be surprisingly competitive for many tasks. If a pretrained encoder
cannot beat a random one, the claim that pretraining is what produces the performance is
unsupported regardless of how good the absolute accuracy is.

\subsection{Preprocessing as a confound}
Physiological noise regression is widely recommended: \citet{fischer2026density} integrate
short-separation regression into cross-validation and report gains of roughly 4\,\% within-subject
and more than 10 percentage points in cross-dataset transfer, and argue that image-space
(parcellated) representations further improve generalisation. Yet the same literature warns that
decoding results can be inflated by non-neural structure: \citet{rybar2024cue} show that cue
presentation windows grossly overestimate semantic BCI performance. Our short-separation ablation
(Section~\ref{sec:results_wss}) speaks directly to this tension.
